% !PRINTED 539
\weberpage{539}
cyklischer Gleichungen. Diesen Satz kann man wieder auf die
Hülfsgleichung $s^{\mathrm{ten}}$ Grades anwenden, und damit so lange fortfahren, bis diese Hülfsgleichung sich auf den ersten Grad reducirt. Damit erhält man also das Resultat:

\gesperrt{Die Lösung einer \gesperrt{Abel}'schen Gleichung lässt sich immer auf die Lösung einer Reihe von cyklischen Gleichungen zurückführen, deren Grade Theiler des Grades der gegebenen Gleichung sind.}

Es ist nicht nothwendig, in die Definition der cyklischen Gleichungen die Irreducibilität mit aufzunehmen. Wir können daher allgemein die Definition so fassen:

\gesperrt{Eine Gleichung $m^{\mathrm{ten}}$ Grades $F(x) = 0$ mit $m$ verschiedenen Wurzeln heisst eine cyklische Gleichung im Körper $\Omega$, wenn ihre Wurzeln $\alpha, \alpha_1, \alpha_2 \ldots \alpha_{m-1}$ nicht rational sind, aber sich so anordnen lassen, dass die cyklischen Functionen der Wurzeln in $\Omega$ rational sind.}

Wenn also
\weq{7}{\pi = (\alpha, \alpha_1, \alpha_2 \ldots \alpha_{m-1})}
ist, so muss jede Function in $\Omega$ enthalten sein, die die Permutationen der cyklischen Gruppe
\weq{8}{C = 1, \pi, \pi^2, \pi^3 \ldots \pi^{m-1}}
gestattet. Die \gesperrt{Galois}'sche Gruppe einer cyklischen Gleichung ist entweder die Periode $C$ selbst und dann ist die Gleichung irreducibel, oder sie ist ein Theiler von $C$; sie besteht dann, wenn $e$ und $f$ zwei ganzzahlige Factoren von $m$ sind und
\[
m = ef
\]
ist, aus den Permutationen
\weq{9}{C_e = 1, \pi^e, \pi^{2e} \ldots \pi^{(f-1)e},}
und die cyklische Gleichung zerfällt in $e$ Factoren $f^{\mathrm{ten}}$ Grades. Denn nehmen wir eine zu der Gruppe $C$ \gesperrt{gehörige} Function $\psi(\alpha, \alpha_1 \ldots \alpha_{m-1})$, so ist diese nach Voraussetzung gleich einer Grösse in $\Omega$, die wir mit $a$ bezeichnen. Auf die rationale Gleichung $\psi = a$ ist dann keine nicht in $C$ enthaltene Permutation anwendbar, und folglich kann die Gruppe der Gleichung keine anderen Substitutionen enthalten als solche, die in $C$ vorkommen. Wenn nun $\pi$ in der Gruppe der Gleichung vorkommt, so ist sie mit $C$ identisch, und da $C$ transitiv ist, ist die Gleichung irreducibel. Ist aber $\pi^e$ die niedrigste Potenz von $\pi$, die in der
%CONT>>
